\begin{abstract}
A hyperedge relates any number of things at once, and many datasets are most naturally hypergraphs. Graph databases store them by reification, which hides each relation's size, members and identity from the storage engine. We describe HStore, an embedded storage engine built for hyperedges. Every structure in HStore, including each hyperedge's membership and the index from an atom to the hyperedges that contain it, is a copy-on-write B+tree over 64-bit keys whose child references carry entry counts and subtree summaries. Small memberships are stored inside the hyperedge's record. Commits log the keys they changed instead of the pages they rewrote, a background checkpoint writes each changed node once, and recovery always rebuilds a prefix of the commit order. Pages are reached through a page directory, so compaction moves them without rewriting any tree, generation or branch. We compare HStore with HyperGraphDB on \M{workloadCount} workloads with asynchronous and synchronous commits and with data larger than both systems' caches. HStore reads incidence sets \M{ratios4asyncreadincidence} to \M{ratios1asyncreadincidence} times faster, writes \M{writtens1asyncwriteupdatehs} per update against \M{writtens1asyncwriteupdatehg}, and ingests hyperedges \M{ratios4asyncingestedges} times faster once the data outgrows the cache. Its median single-commit latency under asynchronous durability is the one measurement where HyperGraphDB is ahead.
\end{abstract}
